\section{What OG--CLEWS Adds to Existing Approaches}
\label{sec:landscape}

CLEWS and OG-Core solve different problems. CLEWS enforces physical balances: electricity
supply meets demand, infrastructure respects capacity, and technologies operate within
resource and policy constraints. OG-Core enforces economic balances among households,
industries, government, saving, investment, and markets. Linking them preserves both
structures while making the relevant economic relationships explicit.

Energy-system and economic models have been combined for decades. MARKAL--MACRO and
MESSAGE--MACRO connect detailed energy planning to aggregate macroeconomic models,
passing energy costs to the economic model and demand back to the energy system
\citep{manne_wene_1992,messner_schrattenholzer_2000}. Their economic representation is
highly aggregated, which limits analysis by age, income, and generation. The related
TIMES/TIAM family provides technology-rich, least-cost energy pathways and can represent
demand response, but it is not itself an aggregate macroeconomic model
\citep{loulou_times_2008}.

Hybrid energy--CGE models incorporate technology detail within, or link it to, a
multi-sector economy \citep{hourcade_hybrid_2006,bohringer_rutherford_2008}. They provide
strong sectoral and price feedback and, when combined with household microdata, can
examine distribution in considerable detail \citep{rausch_2011}. The cited applications
do not use the forward-looking cohort structure of an overlapping-generations model.
OLG climate models do address burdens across generations. The model of
\citet{kotlikoff_2021}, for example, includes coal, oil, gas, clean energy, extraction
costs, and climate damages, but not a technology- and infrastructure-resolved national
resource plan.

IEEM and CLEWs--IO are the closest comparisons on the resource side. IEEM combines a
dynamic CGE model with environmental accounts, spatial land-use change, ecosystem-service
modelling, and endogenous feedback between natural capital and the economy
\citep{banerjee_ieem_2020}. CLEWs--IO soft-links the resource model to a regional
input--output framework, estimates direct, indirect, and induced employment, and allows
those employment effects to reshape the resource pathway
\citep{bararzadeh_clews_io_2026}. OG--CLEWS extends this resource--economy literature in a
different direction: it connects CLEWS to an economy in which households make
forward-looking decisions over their lives, generations overlap, policy changes affect
the government's budget over time, and changes in health can alter survival, labour, and
population.

\begin{table}[H]
\centering\scriptsize
\caption{Comparison with related resource--economy approaches.}
\label{tab:comparison}
\begin{tabularx}{\linewidth}{@{}>{\bfseries\raggedright\arraybackslash}p{0.19\linewidth} >{\raggedright\arraybackslash}p{0.30\linewidth} L@{}}
\toprule
Approach & Main emphasis & What OG--CLEWS adds \\
\midrule
Classic energy--macro & energy technologies linked to aggregate growth and demand & distribution by age and lifetime income, overlapping generations, and explicit fiscal incidence \\
TIMES/TIAM & technology-rich least-cost energy pathways and demand response & household, industry, demographic, and fiscal consequences in a general-equilibrium economy \\
Hybrid energy--CGE & technology, sectors, prices, and household groups & forward-looking life-cycle choices, generational incidence, and an age-structured population \\
OLG climate models & fuels, clean energy, climate damages, and intergenerational policy & a detailed national resource plan with explicit technologies, infrastructure, system costs, and emissions \\
IEEM & environmental accounts, spatial land use, ecosystem services, and natural-capital feedback & household life-cycle behaviour, generational distribution, and intertemporal public finance \\
CLEWs--IO & direct, indirect, and induced employment linked to CLEWS resource choices & behavioural responses in consumption, production, saving, investment, taxation, transfers, and public debt \\
OG--CLEWS & detailed resource choices and an overlapping-generations economy & a single assessment of technological, sectoral, household, demographic, health, and fiscal consequences \\
\bottomrule
\end{tabularx}
\end{table}

The advance is the combination of capabilities. OG--CLEWS retains the technological and
physical detail of a country CLEWS model while using a country OG-Core calibration to
trace the consequences through industries, households of different ages and lifetime
incomes, the population, and the public budget. It can therefore distinguish effects that
aggregate models combine: a change in system cost from a change in taxation, public
infrastructure from private generation capital, and a near-term labour effect from a
longer-term demographic effect. Economic activity and the equilibrium return can then
inform the resource plan's demand and financing assumptions. This provides a coherent
account of who is affected, through which mechanism, and when.
